%----------------------------------------------------------------------------------------
% GEVOELIGHEID EN GROTERE PROBLEMEN. Alle opgaven zijn eigen werk.
%
% Opbouw: Probleem, Start, leerstof, en op een nieuwe pagina de Opdrachten, Koppel en Test jezelf.
% Labels van dit hoofdstuk beginnen met h3-.
%----------------------------------------------------------------------------------------

\mychapter{3}{Gevoeligheid en grotere problemen}

\begin{probleem}
\vspace{-1mm}\naast{De mini-onderneming uit hoofdstuk 2 maakt 135 kaarsen en 30 dozen zeep, voor 690 euro winst. Dan komen er drie voorstellen. Een winkel wil de kaarsen voor meer verkopen: tot welke winst per kaars blijft dit het beste plan? Een leverancier biedt extra geurolie aan, voor 8 cent per ml: loont dat, en hoeveel koop je? En een grootvader wil een paar uur komen helpen, als jullie hem iets betalen: hoeveel is een extra uur werk waard?}{\begin{tikzpicture}[x=.02cm,y=.02cm,font=\scriptsize]
  \path[gebied] (0,0) -- (150,0) -- (135,30) -- (30,100) -- (0,100) -- cycle;
  \draw[grens] (0,0) -- (150,0) -- (135,30) -- (30,100) -- (0,100) -- cycle;
  \draw[grijs,densely dashed] (152.5,0) -- (130,45) (60,100) -- (165,30);
  \optimum[below left]{135}{30}{}
\end{tikzpicture}}
\end{probleem}

\lkr{Leerplan Complementair (III-WiCo-d): LPD 39 (sectie 1 en 2) en LPD 38 (sectie 3), met de wenken over het bereik van de coëfficiënten van de doelfunctie, de invloed van de beperkingen, de schaduwprijs en het nut van een gevoeligheidsanalyse, en voor ICT over het opstellen van het model en problemen met meer dan twee variabelen. Voorkennis: hoofdstuk 2, vectoren (een basis van twee niet-evenwijdige vectoren) en stelsels van twee vergelijkingen. In sectie 3 werken de leerlingen met een rekenblad met een oplosser: plan daarvoor een les in een computerklas of geef het als taak voor thuis. Het boek toont een rekenblad op papier, zodat de les ook zonder computer kan.}

In hoofdstuk 2 vond je het beste plan bij gegeven getallen. Maar die getallen liggen zelden vast. Een prijs verandert, een leverancier levert meer of minder, en een winst per stuk is vaak maar een schatting. Moet je dan alles opnieuw uitrekenen? Of kun je zeggen hoe ver een getal mag veranderen voor het plan verandert? Dat onderzoek heet een \emph{gevoeligheidsanalyse}.

In dit hoofdstuk leer je twee dingen. Eerst hoe gevoelig het optimum is voor de doelfunctie en voor de beperkingen, met de schaduwprijs als belangrijkste getal. Daarna hoe je een probleem met meer dan twee onbekenden oplost: dat kan niet meer met een tekening, wel met een computer.

\begin{oefeningen}[Start]
\oef{h3-start-1}

\oef{h3-verken}
\end{oefeningen}

\lkr{Laat oefening \ref{oef:h3-verken} maken vóór de uitleg. De leerlingen zien dat het optimum blijft liggen als de winst een beetje verandert, en pas springt als de niveaulijn evenwijdig wordt met een zijde.}

%========================================================================================
\section{De doelfunctie verandert}

\subsection{Wanneer blijft het optimum liggen?}

Neem het fietsatelier uit oefening~\ref{oef:h2-model}: de winst is $W = 30x + 50y$, en het beste plan is 20 stadsfietsen en 10 racefietsen. Wat als een stadsfiets 35 euro opbrengt in plaats van 30? De niveaulijnen kantelen een beetje. Zolang ze het gebied als laatste in hetzelfde hoekpunt raken, blijft het plan hetzelfde. Alleen de winst verandert.

Om te weten hoe ver je mag gaan, hoef je niet alle hoekpunten te vergelijken. De volgende stelling zegt dat het volstaat om een hoekpunt met zijn twee buren te vergelijken, de hoekpunten aan het andere eind van de twee zijden door dat hoekpunt.

\begin{theorie}
Een beperkende voorwaarde is \textbf{\textcolor{accent}{bindend}} in een punt als daar de gelijkheid geldt: het punt ligt op haar grensrechte. Het verschil tussen het rechterlid en het linkerlid heet de \textbf{\textcolor{accent}{speling}}.

\begin{minipage}[t]{.64\linewidth}
\vspace{0pt}
\begin{stelling}{Buurhoekpunten}
Neem een begrensd oplossingsgebied $G$, een hoekpunt $P$ en zijn twee buren $A$ en $B$. Is

\vspace{1mm}
\qquad $\formule{Z(A) \leqslant Z(P)}$ \ en \ $\formule{Z(B) \leqslant Z(P)}$,

\vspace{1mm}
dan is $Z$ maximaal in $P$.
\end{stelling}
\end{minipage}\hfill
\begin{minipage}[t]{.32\linewidth}
\vspace{0pt}
\centering
\begin{tikzpicture}[scale=.62,font=\scriptsize]
  \fill[accent!10] (3,1) -- (-1.5,1) -- (-1.5,4.2) -- (2.2,4.2) -- cycle;
  \path[gebied,draw=grijs,thick] (3,1) -- (0,1) -- (-.6,2.6) -- (.8,3.6) -- (2.6,2.6) -- cycle;
  \draw[->,accent,thick] (3,1) -- (0,1) node[below,midway,text=black] {$\vec u$};
  \draw[->,accent,thick] (3,1) -- (2.6,2.6) node[right,midway,text=black] {$\vec v$};
  \fill (3,1) circle (1.8pt) node[right] {$P$};
  \fill (0,1) circle (1.6pt) node[below left] {$A$};
  \fill (2.6,2.6) circle (1.6pt) node[above right] {$B$};
\end{tikzpicture}
\end{minipage}

\begin{bewijs}
Stel $\vec u = A - P$ en $\vec v = B - P$. Ze zijn niet evenwijdig, dus elk punt $Q$ is te schrijven als $Q = P + s\,\vec u + t\,\vec v$. De zijde $[PA]$ ligt op een grensrechte $f(x, y) = 0$, met $f$ van de eerste graad en $G$ aan de kant waar $f \geqslant 0$. Dan is $f(P) = f(A) = 0$, en
\bewijsregel{f(Q) = f(P) + s\,\bigl(f(A) - f(P)\bigr) + t\,\bigl(f(B) - f(P)\bigr) = t\,f(B) .}
Omdat $f(Q) \geqslant 0$ en $f(B) > 0$, is $t \geqslant 0$. Net zo is $s \geqslant 0$, met de grensrechte van $[PB]$. Op dezelfde manier is
\bewijsregel{Z(Q) = Z(P) + s\,\bigl(Z(A) - Z(P)\bigr) + t\,\bigl(Z(B) - Z(P)\bigr) \leqslant Z(P) .}
Dat geldt voor elk punt $Q$ in de twee halfvlakken van de bindende voorwaarden in $P$ (de getinte hoek in de figuur), en dus zeker voor elk punt van $G$.
\end{bewijs}
\end{theorie}

In woorden: wie in een hoekpunt staat en bij geen van beide buren beter wordt, staat in het optimum. Een plaatselijk beste hoekpunt is dus ook het beste van allemaal. Dat komt weer doordat het gebied convex is. Het bewijs toont zelfs meer: in de hele getinte hoek is $Z$ nergens groter dan in $P$. Dat gebruiken we in sectie 2.

\begin{theorie}
Het \textbf{\textcolor{accent}{optimaliteitsbereik}} van een coëfficiënt van de doelfunctie is het interval waarin die coëfficiënt mag liggen, terwijl de andere vast blijven, zonder dat het optimale hoekpunt verandert.
\end{theorie}

\voorbeeld{}{Het fietsatelier heeft het toegestane gebied met hoekpunten $(0, 0)$, $(25, 0)$, $(20, 10)$, $(10, 15)$ en $(0, 15)$. Tussen welke grenzen mag de winst $c$ op een stadsfiets liggen, als 20 stadsfietsen en 10 racefietsen het beste plan moet blijven? De winst op een racefiets blijft 50 euro.\label{vb:h3-bereik}}{%
\uitleg{De buren van $P(20, 10)$ zijn $A(25, 0)$ en $B(10, 15)$. Volgens de stelling volstaat het dat $P$ niet slechter is dan die twee.}
\naast{\begin{stappenvrij}{5.2cm}{3.4cm}
W = cx + 50y & De doelfunctie met de onbekende winst $c$. \\[1ex]
20c + 500 \geqslant 25c & $W(P) \geqslant W(A)$. \\[.5ex]
\Rightarrow \ c \leqslant 100 & \\[1ex]
20c + 500 \geqslant 10c + 750 & $W(P) \geqslant W(B)$. \\[.5ex]
\Rightarrow \ c \geqslant 25 & \\[1ex]
25 \leqslant c \leqslant 100 & Het optimaliteitsbereik.
\end{stappenvrij}}{\begin{tikzpicture}[scale=.12,font=\scriptsize]
  \path[gebied] (0,0) -- (25,0) -- (20,10) -- (10,15) -- (0,15) -- cycle;
  \roosterom{0}{5}{10}{40}{0}{5}{10}{25}
  \draw[grens] (0,0) -- (25,0) -- (20,10) -- (10,15) -- (0,15) -- cycle;
  \begin{scope}\venster{0}{41}{0}{26}\grenslijn[niveau]{25}{50}{1000}\grenslijn[niveau]{100}{50}{2500}\end{scope}
  \optimum[above right]{20}{10}{$P$}
  \hoekpunt[below]{25}{0}{$A$}
  \hoekpunt[above]{10}{15}{$B$}
\end{tikzpicture}}
Bij $c = 25$ is de niveaulijn evenwijdig met de zijde $[BP]$, bij $c = 100$ met $[PA]$. Daartussen blijft $P$ optimaal. Bij een winst boven 100 euro herstelt het atelier best alleen stadsfietsen, onder 25 euro schuift het optimum naar $B$.}

Er is een tweede manier, met de richtingscoëfficiënten. Het hoekpunt blijft optimaal zolang de niveaulijnen niet steiler zijn dan de ene zijde door $P$ en niet vlakker dan de andere. Dat is vaak sneller, maar let op de tekens.

\voorbeeld{}{Bepaal met de richtingscoëfficiënten het optimaliteitsbereik van de winst $q$ op een racefiets, als de winst op een stadsfiets 30 euro blijft.\label{vb:h3-rico}}{%
\uitleg{De bindende voorwaarden in $P(20, 10)$ zijn $x + 2y \leqslant 40$ (werktijd) en $2x + y \leqslant 50$ (onderdelen). De rico van de niveaulijnen $30x + qy = k$ moet tussen hun rico's liggen.}
\begin{stappen}
-2 \leqslant -\dfrac{30}{q} \leqslant -\dfrac12 & De rico's van $2x + y = 50$ en $x + 2y = 40$. \\[3ex]
\dfrac12 \leqslant \dfrac{30}{q} \leqslant 2 & Vermenigvuldig met $-1$: de tekens keren om. \\[3ex]
15 \leqslant q \leqslant 60 & Voor $q > 0$ omkeren: $\tfrac{q}{30}$ ligt tussen $\tfrac12$ en 2.
\end{stappen}
Controle met de buren: $600 + 10q \geqslant 750$ geeft $q \geqslant 15$, en $600 + 10q \geqslant 300 + 15q$ geeft $q \leqslant 60$.}

\begin{werkwijze}
Vervang de coëfficiënt door een letter. Zoek de twee buren van het optimale hoekpunt $P$, en eis dat $Z(P)$ minstens zo groot is als $Z$ in elke buur (voor een minimum: hoogstens zo groot). Los de twee ongelijkheden op. Controleer met de richtingscoëfficiënten als dat kan.
\end{werkwijze}

\nujij{Wat is het optimaliteitsbereik van de winst op een stadsfiets als een racefiets 40 euro opbrengt?}{$20c + 400 \geqslant 25c$ en $20c + 400 \geqslant 10c + 600$: $20 \leqslant c \leqslant 80$.}

\begin{oefeningen}
\oef{h3-buur}

\oef{h3-bereik}
\end{oefeningen}

%========================================================================================
\section{Een beperking verandert}

\subsection{De schaduwprijs}

Het fietsatelier kan een extra uur werk inhuren. Hoeveel mag dat uur kosten? Met een extra uur verschuift de grensrechte $x + 2y = 40$ een beetje naar buiten, evenwijdig met zichzelf. Het optimale hoekpunt schuift mee langs de andere bindende grensrechte, en de winst stijgt. Die stijging is het getal dat we zoeken.

\begin{theorie}
De \textbf{\textcolor{accent}{schaduwprijs}} van een beperkende voorwaarde is de verandering van de optimale waarde van de doelfunctie als het rechterlid van die voorwaarde met 1 toeneemt.

\begin{stelling}{Schaduwprijs}
Het optimale hoekpunt $P$ is het snijpunt van de grensrechten van twee bindende voorwaarden. Verandert het rechterlid van een ervan met $h$, dan schuift $P$ naar $P_h = P + h\,\vec w$, met een vaste vector $\vec w$. Zolang $P_h$ aan alle voorwaarden voldoet, is $P_h$ optimaal en

\vspace{1mm}
\qquad $\formule{Z(P_h) = Z(P) + h \cdot s}$ \qquad met een vaste schaduwprijs $\formule{s = Z(\vec w)}$.
\end{stelling}

\begin{bewijs}
Stel dat $P$ het stelsel $a_1x + b_1y = c_1$, $a_2x + b_2y = c_2$ oplost, met $D = a_1b_2 - a_2b_1 \neq 0$. Vervang $c_1$ door $c_1 + h$. Dan is
\bewijsregel{x = \dfrac{(c_1 + h)\,b_2 - c_2b_1}{D}, \qquad y = \dfrac{a_1c_2 - a_2(c_1 + h)}{D} ,}
dus $P_h = P + h\,\vec w$ met $\vec w = \tfrac1D(b_2, -a_2)$. Omdat $Z$ van de eerste graad is, is $Z(P_h) = Z(P) + h\,Z(\vec w)$. De twee grensrechten schuiven evenwijdig, dus de zijden door $P_h$ hebben dezelfde richtingen $\vec u$ en $\vec v$ als die door $P$. In het bewijs van de buurhoekpunten tellen alleen die richtingen: $Z(\vec u) \leqslant 0$ en $Z(\vec v) \leqslant 0$ blijven gelden, dus $P_h$ blijft optimaal.
\end{bewijs}
\end{theorie}

De schaduwprijs zegt dus hoeveel één extra eenheid van een middel waard is. Hij blijft dezelfde zolang hetzelfde paar grensrechten het optimum bepaalt. Dat interval van het rechterlid heet het \emph{geldigheidsbereik}. Daarbuiten springt het optimum naar een ander hoekpunt, met een andere schaduwprijs.

\voorbeeld{}{Bereken de schaduwprijs van de werktijd van het fietsatelier ($x + 2y \leqslant 40$, in uren) en het geldigheidsbereik.\label{vb:h3-schaduw}}{%
\uitleg{We vervangen 40 door een letter $T$, lossen het stelsel van de twee bindende grensrechten op, en schrijven de winst als functie van $T$.}
\naast{\begin{stappenvrij}{5cm}{3.6cm}
\begin{cases} x + 2y = T \\ 2x + y = 50 \end{cases} & De bindende grensrechten. \\[3ex]
x = \tfrac{100 - T}{3}, \ y = \tfrac{2T - 50}{3} & Oplossen. Controle: $T = 40$ geeft $(20, 10)$. \\[1ex]
W = \tfrac{500 + 70T}{3} & $30x + 50y$, uitgewerkt. \\[1ex]
s = \tfrac{70}{3} \approx 23{,}33 & Euro per extra uur werk. \\[1ex]
25 \leqslant T \leqslant 47{,}5 & Uit $y \geqslant 0$ en $y \leqslant 15$.
\end{stappenvrij}}{\begin{tikzpicture}[scale=.12,font=\scriptsize]
  \path[gebied] (0,0) -- (25,0) -- (20,10) -- (10,15) -- (0,15) -- cycle;
  \roosterom{0}{5}{10}{40}{0}{5}{10}{25}
  \draw[grens] (0,0) -- (25,0) -- (20,10) -- (10,15) -- (0,15) -- cycle;
  \begin{scope}\venster{0}{41}{0}{26}\grenslijn[densely dashed]{1}{2}{44}\grenslijn{2}{1}{50}\end{scope}
  \optimum[below left]{20}{10}{$P$}
  \fill[accent] (56/3,38/3) circle (2.4pt) node[above right,text=black] {$P_4$};
\end{tikzpicture}}
Elk extra uur werk levert dus $23{,}33$ euro meer winst op, tot 47,5 uur. Dan zit het atelier aan de 15 racefietsen, en heeft een extra uur minder nut. Een uur inhuren loont als het minder dan 23,33 euro kost.}

En de racefietsen? In $P$ worden er 10 hersteld, terwijl er 15 binnenkomen: de voorwaarde $y \leqslant 15$ is niet bindend, met speling 5. Komt er een racefiets meer binnen, dan verandert er niets.

\begin{theorie}
\begin{stelling}{Niet-bindende voorwaarde}
Een voorwaarde die in het optimale hoekpunt niet bindend is, heeft schaduwprijs 0, zolang haar speling niet negatief wordt.
\end{stelling}

\begin{bewijs}
Verander het rechterlid van zo'n voorwaarde zo weinig dat $P$ er nog aan voldoet. De bindende voorwaarden veranderen niet, dus het nieuwe gebied ligt nog altijd in de getinte hoek van het bewijs van de buurhoekpunten, waar $Z \leqslant Z(P)$. En $P$ ligt in het nieuwe gebied. Dus blijft $P$ optimaal, met dezelfde waarde.
\end{bewijs}
\end{theorie}

\begin{werkwijze}
Schrijf de twee bindende voorwaarden met een letter als rechterlid. Los het stelsel op en vul in in de doelfunctie: de coëfficiënt van de letter is de schaduwprijs. Het geldigheidsbereik vind je door te eisen dat het hoekpunt aan de andere voorwaarden voldoet. Een voorwaarde die niet bindend is, heeft schaduwprijs 0.
\end{werkwijze}

\nujij{Wat is de schaduwprijs van het onderdelenbudget van het fietsatelier, $20x + 10y \leqslant 500$, in euro winst per euro budget?}{Met $x + 2y = 40$ en $20x + 10y = B$: $y = \tfrac{800 - B}{30}$, $W = 30(40 - 2y) + 50y = 1200 - 10y = \tfrac{B}{3} + \tfrac{1000}{3}$. De schaduwprijs is $\tfrac13$ euro per euro.}

\begin{oefeningen}
\oef{h3-speling}

\oef{h3-schaduw}
\end{oefeningen}

\subsection{Terug naar het probleem}

\voorbeeld{}{Terug naar het probleem. Beantwoord de drie vragen van de mini-onderneming: het optimaliteitsbereik van de winst per kaars, de waarde van extra olie, en de waarde van een extra uur werk.\label{vb:h3-terug}}{%
\uitleg{Het optimum is $P(135, 30)$, met de buren $A(150, 0)$ en $B(30, 100)$. Werktijd en olie zijn bindend, de verpakkingen niet.}
\begin{stappen}
135c + 150 \geqslant 150c, \ 135c + 150 \geqslant 30c + 500 & Winst $c$ per kaars: $W(P)$ tegen de buren. \\[1ex]
\tfrac{10}{3} \leqslant c \leqslant 10 & Tussen 3,33 en 10 euro per kaars blijft het plan. \\[1ex]
\begin{cases} 20x + 30y = 3600 \\ 10x + 5y = b \end{cases} & Olie: $b$ ml. \\[3ex]
y = \tfrac{1800 - b}{10}, \ W = 540 + 0{,}1\,b & Geldig voor $800 \leqslant b \leqslant 1800$ (uit $0 \leqslant y \leqslant 100$). \\[1ex]
\begin{cases} 20x + 30y = T \\ 10x + 5y = 1500 \end{cases} & Werktijd: $T$ minuten. \\[3ex]
y = \tfrac{T - 3000}{20}, \ W = 150 + 0{,}15\,T & Geldig voor $3000 \leqslant T \leqslant 5000$.
\end{stappen}
Een kaars mag tot 10 euro opbrengen zonder dat het plan verandert. Een ml olie is 10 cent waard: aan 8 cent loont het, en je koopt 300 ml bij (tot 1800 ml), voor $300 \cdot 0{,}02 = 6$ euro extra. Een minuut werk is 15 cent waard, dus 9 euro per uur, tot ruim 83 uur. Grootvader mag dus hoogstens 9 euro per uur krijgen, wil het voor de onderneming lonen. Extra verpakkingen hebben geen waarde.}

Waarom is zo'n analyse nuttig? De getallen in een model zijn vaak schattingen: de winst per stuk, de tijd per stuk. De optimaliteitsbereiken tonen hoe zeker het plan is: een ruim bereik betekent dat een kleine schattingsfout niets uitmaakt. En de schaduwprijzen tonen waar je best investeert: in het middel met de hoogste schaduwprijs, en nooit meer dan die prijs.

\begin{oefeningen}
\oef{h3-terug}
\end{oefeningen}

%========================================================================================
\section{Grotere problemen met ICT}

\subsection{Meer dan twee onbekenden}

De mini-onderneming wil ook geurzakjes maken: 10 minuten werk en 5 ml olie per zakje, voor 3 euro winst. Nu zijn er drie onbekenden, $x$, $y$ en $z$. Het toegestane gebied ligt in de ruimte: elke voorwaarde geeft een halfruimte, en samen vormen ze een veelvlak. Hoekpunten, buren en de hoekpuntstelling bestaan nog altijd, maar tekenen wordt lastig, en bij tien of honderd onbekenden onmogelijk.

\begin{theorie}
Een \textbf{\textcolor{accent}{lineair programmeringsprobleem}} met $n$ onbekenden $x_1, \ldots, x_n \geqslant 0$ vraagt het maximum of het minimum van

\vspace{1mm}
\qquad $\formule{Z = c_1x_1 + c_2x_2 + \cdots + c_nx_n}$

\vspace{1mm}
onder beperkende voorwaarden van de vorm $\formule{a_1x_1 + a_2x_2 + \cdots + a_nx_n \leqslant b}$ (of $\geqslant$, of $=$).
\end{theorie}

\naast{Computers lossen zulke problemen op met de \emph{simplexmethode}, die George Dantzig in 1947 bedacht. Het idee volgt uit de stelling van de buurhoekpunten: begin in een hoekpunt, ga naar een buur waar $Z$ groter is, en herhaal dat. Is geen enkele buur beter, dan zit je in het optimum. In de praktijk vindt de methode zo het optimum tussen miljoenen hoekpunten, na relatief weinig stappen. De stelling van de buurhoekpunten bewezen we alleen in het vlak; in de ruimte en in meer dimensies geldt ze ook.}{\begin{tikzpicture}[x={(-.009cm,-.005cm)},y={(.024cm,0cm)},z={(0cm,.011cm)},font=\scriptsize]
  \fill[accent!12] (0,0,300) -- (0,30,270) -- (135,30,0) -- (150,0,0) -- cycle;
  \draw[grijs,->] (0,0,0) -- (185,0,0) node[below] {$x$};
  \draw[grijs,->] (0,0,0) -- (0,125,0) node[right] {$y$};
  \draw[grijs,->] (0,0,0) -- (0,0,340) node[left] {$z$};
  \foreach \a/\b in {{(0,0,0)}/{(0,0,300)},{(0,0,0)}/{(0,100,0)},{(0,0,0)}/{(150,0,0)},{(0,0,300)}/{(0,30,270)},{(0,0,300)}/{(150,0,0)},{(0,100,0)}/{(0,100,60)},{(0,100,0)}/{(30,100,0)},{(0,30,270)}/{(0,100,60)},{(0,30,270)}/{(135,30,0)},{(0,100,60)}/{(30,100,0)},{(150,0,0)}/{(135,30,0)},{(135,30,0)}/{(30,100,0)}}
    \draw[accent!70!black,thick] \a -- \b;
  \fill[accent] (0,30,270) circle (2.6pt);
  \foreach \p in {(0,0,300),(0,100,0),(0,100,60),(150,0,0),(135,30,0),(30,100,0)} \fill \p circle (1.4pt);
\end{tikzpicture}}

De figuur toont het toegestane gebied van de mini-onderneming met geurzakjes: een veelvlak met acht hoekpunten. Het grote punt is het optimum, dat we hieronder met een rekenblad vinden.

\subsection{Met een rekenblad}

De meeste rekenbladen hebben een \textit{oplosser} (Excel en LibreOffice Calc noemen hem Oplosser of Solver) die lineaire programmeringsproblemen oplost. Het werk zit niet in het rekenen, maar in het model: welke cellen veranderen, welke cel moet zo groot mogelijk worden, en welke voorwaarden gelden.

\voorbeeld{}{De mini-onderneming maakt $x$ kaarsen, $y$ dozen zeep en $z$ geurzakjes. Een zakje vraagt 10 minuten werk en 5 ml olie, en brengt 3 euro winst op. Zet het model in een rekenblad en laat de oplosser het beste plan zoeken.\label{vb:h3-rekenblad}}{%
\uitleg{We zetten de gegevens in een tabel: per product een kolom, per voorwaarde een rij. Zo zie je meteen of je alles hebt, en de formules zijn telkens dezelfde.}
\begin{stappen}
W = 4x + 5y + 3z & De doelfunctie. \\[1ex]
20x + 30y + 10z \leqslant 3600 & Werktijd in minuten. \\
10x + 5y + 5z \leqslant 1500 & Olie in ml. \\
y \leqslant 100, \quad x, y, z \geqslant 0 & Verpakkingen.
\end{stappen}
{\centering\small\renewcommand{\arraystretch}{1.3}\arrayrulecolor{grijs}%
\begin{tabular}{|c|l|c|c|c|l|c|}
\hline
 & \textbf{A} & \textbf{B} & \textbf{C} & \textbf{D} & \textbf{E} & \textbf{F} \\ \hline
1 & & kaars & zeep & zakje & gebruikt & beschikbaar \\ \hline
2 & aantal & \cellcolor{accent!15}0 & \cellcolor{accent!15}30 & \cellcolor{accent!15}270 & & \\ \hline
3 & winst & 4 & 5 & 3 & \cellcolor{accent!30}\texttt{=SOMPRODUCT(B2:D2;B3:D3)} & \\ \hline
4 & werk & 20 & 30 & 10 & \texttt{=SOMPRODUCT(B\$2:D\$2;B4:D4)} & 3600 \\ \hline
5 & olie & 10 & 5 & 5 & \texttt{=SOMPRODUCT(B\$2:D\$2;B5:D5)} & 1500 \\ \hline
6 & verpakking & 0 & 1 & 0 & \texttt{=SOMPRODUCT(B\$2:D\$2;B6:D6)} & 100 \\ \hline
\end{tabular}\par}
\medskip
In de oplosser: doel E3 maximaliseren, door B2:D2 te veranderen, met de voorwaarden E4:E6 $\leqslant$ F4:F6, variabelen niet negatief, en de methode voor lineaire problemen (simplex). De oplosser geeft $x = 0$, $y = 30$ en $z = 270$, voor 960 euro.

Interpretatie: met geurzakjes stoppen de kaarsen helemaal. Een zakje vraagt per euro winst veel minder werk dan een kaars. Werk en olie zijn opnieuw volledig op.}

\begin{werkwijze}
Zet in een rekenblad een rij met de onbekenden (eerst willekeurige getallen), een tabel met de coëfficiënten, een cel met de doelfunctie (\texttt{SOMPRODUCT}) en per voorwaarde een cel met het linkerlid en een met het rechterlid. Vul in de oplosser het doel, de veranderlijke cellen en de voorwaarden in, en kies niet-negatief en lineair. Controleer het resultaat: voldoet het aan alle voorwaarden, en is het zinvol in de context?
\end{werkwijze}

\subsection{Het gevoeligheidsrapport}

Veel oplossers maken op vraag een \emph{gevoeligheidsrapport}. Daarin staan per voorwaarde de schaduwprijs en het geldigheidsbereik, en per onbekende het optimaliteitsbereik van haar coëfficiënt. Bij drie onbekenden kun je die getallen niet meer uit een tekening halen, maar ze betekenen hetzelfde als in sectie 2.

\voorbeeld{}{Het rapport geeft voor de mini-onderneming met geurzakjes deze schaduwprijzen: werktijd 0,10 euro per minuut, olie 0,40 euro per ml, verpakkingen 0. Waarom maakt de onderneming geen kaarsen meer? Hoeveel winst moet een kaars opleveren om weer mee te doen?\label{vb:h3-rapport}}{%
\uitleg{Een kaars gebruikt werk en olie die nu naar zeep en zakjes gaan. Met de schaduwprijzen rekenen we uit wat die middelen elders waard zijn.}
\begin{stappen}
20 \cdot 0{,}10 + 10 \cdot 0{,}40 = 6 & Wat het werk en de olie van één kaars elders opbrengen, in euro. \\[1ex]
4 < 6 & Een kaars brengt minder op dan ze kost aan gemiste winst. \\[1ex]
c > 6 & Pas boven 6 euro winst per kaars loont het om kaarsen te maken.
\end{stappen}
Ter controle: met $c = 7$ geeft de oplosser weer 135 kaarsen, 30 dozen zeep en geen zakjes, voor 1095 euro. Bij precies 6 euro zijn beide plannen even goed: 960 euro.}

Zo zijn schaduwprijzen ook prijzen voor beslissingen: een nieuw product is pas interessant als het meer opbrengt dan de schaduwprijzen van alles wat het verbruikt.

\begin{oefeningen}
\oef{h3-model}

\oef{h3-rapport}
\end{oefeningen}

%========================================================================================
\opdrachtenkop

\begin{oefeningen}
\oef{h3-o1}

\oef{h3-o2}

\oef{h3-o3}

\oef{h3-o4}

\oef{h3-o5}

\oef{h3-o6}

\oef{h3-o7}

\oef{h3-o8}

\oef{h3-o9}

\oef{h3-o10}

\oef{h3-o11}

\oef{h3-o12}

\oef{h3-o13}

\oef{h3-o14}

\oef{h3-o15}

\oef{h3-o16}

\oef{h3-o21}

\oef{h3-o17}

\oef{h3-o18}

\oef{h3-o19}

\oef{h3-o20}
\end{oefeningen}

%----------------------------------------------------------------------------------------
% KOPPEL: verbanden met andere leerstof
%----------------------------------------------------------------------------------------
\needspace{20\baselineskip}
\opdrachtsectie{Koppel}

\lkr{Afgeleiden en Matrices zagen de leerlingen in het 5de jaar. Ruimtemeetkunde staat in het 6de jaar: geef die oefening pas als de snijpunten van vlakken gezien zijn.}

\begin{oefeningen}
\oef{h3-koppel-1}

\oef{h3-koppel-2}

\oef{h3-koppel-3}
\end{oefeningen}

%----------------------------------------------------------------------------------------
\needspace{30\baselineskip} % de test en het schema Fout? Herhaal samen op een pagina
\opdrachtsectie{Test jezelf}

\begin{oefeningen}
\oef{h3-t1}

\oef{h3-t2}

\oef{h3-t3}

\oef{h3-t4}

\oef{h3-t5}

\oef{h3-t6}
\end{oefeningen}

\herhaalschema{%
  \fv{oef:h3-t1}{vb:h3-bereik}{oef:h3-bereik}%
  \fv{oef:h3-t2}{vb:h3-rico}{oef:h3-bereik}%
  \fv{oef:h3-t3}{vb:h3-schaduw}{oef:h3-schaduw}%
  \fv{oef:h3-t4}{vb:h3-terug}{oef:h3-terug}%
  \fv{oef:h3-t5}{vb:h3-rekenblad}{oef:h3-model}%
  \fv{oef:h3-t6}{vb:h3-rapport}{oef:h3-rapport}}

\ifzelfstudie
\section*{Maak je eigen samenvatting}
Sluit dit boek en schrijf in je schrift: wat bindend en speling betekenen, de stelling van de buurhoekpunten en het idee van haar bewijs, het optimaliteitsbereik en de twee manieren om het te berekenen, de schaduwprijs en waarom hij constant is binnen het geldigheidsbereik, de schaduwprijs van een niet-bindende voorwaarde, waarom een gevoeligheidsanalyse nuttig is, hoe je een probleem met meer onbekenden in een rekenblad zet, en wat je met de schaduwprijzen uit een rapport doet. Vergelijk daarna met de kaders.
\fi
